\chapter{Introduzione e contesto}
\label{cap1}
\section{Introduzione e motivazioni}
Nelle proprie interazioni, l'essere umano ha la capacità di anticipare il comportamento degli altri individui e di agire di conseguenza. Questa abilità di anticipare un'azione è intrinseca della nostra specie ed è involontaria, e ne facciamo largo uso in quasi tutte le interazioni umane. Lo studio di questa capacità all'interno del campo della \textit{computer vision} prende il nome di Human Activity Analysis (HAA), la quale presenta una vasta gamma di applicazioni, tra le quali figurano il rilevamento di comportamenti anomali nell'ambito della videosorveglianza, l'avviso di pericolo nei sistemi avanzati di assistenza alla guida, le interazioni uomo-macchina, la realtà aumentata e il recupero video. È chiaro come il comune denominatore di questi svariati impieghi risieda nella comprensione video, ovvero ricavare informazioni dall'osservazione delle sequenze di azioni. 

\section{Stato dell'arte}
Per comprendere la specificità dell'Action Anticipation nell'ampio panorama della comprensione video occorre prima di tutto distinguerla da Action Recognition e Action Prediction (o Early Action Recognition). Per Action Recognition si intende la classificazione di un'azione umana basandosi sulla sua completa esecuzione.
L'Action Prediction, che spesso viene chiamata anche Early Recognition, mira a riconoscere un'azione attualmente in corso, a partire da un'osservazione parziale. Il suo obiettivo è riconoscere l'azione il prima possibile, mentre è ancora in fase di svolgimento. 
Per Action Anticipation, invece, si intende la previsione di un'azione ancor prima che essa abbia avuto inizio. L'osservazione termina prima che l'azione futura abbia inizio. Questa distinzione non è banale in quanto, nell'action recognition, il modello ha accesso tutti i componenti che caratterizzano l'azione: la fase iniziale, lo svolgimento centrale e la fase conclusiva, mentre nell'early recognition ha accesso almeno alla fase iniziale, contrariamente allo scenario dell'anticipation, dove il modello deve prevedere un evento completamente nuovo sfruttando una comprensione profonda di segnali indiretti, ovvero le azioni precedenti, gli oggetti presenti nella scena, la postura del corpo, la direzione dello sguardo, tutti elementi che costituiscono il contesto da cui dedurre l'intenzione futura. Questa differenza ha profonde implicazioni sulle differenze di difficoltà dei task: l'anticipation è intrinsecamente più incerta della recognition poiché, dall'osservazione di uno stesso evento si ricavano più azioni future plausibili. 

\section{Obiettivo e contributo del lavoro}
